% Documento da pesquisa: Gestão Hospitalar no SUS
% Partes atuais: Análise Exploratória (seções 1 a 9) e Estimação
% (seções 10 a 13) do painel de hospitais SUS/SP.
% Linguagem adaptada para leitura por gestores (não estatísticos).
% Compilar no Overleaf com pdfLaTeX. Ver LEIAME_overleaf.txt.
% Este arquivo substitui o antigo main.tex.

\documentclass[12pt,a4paper]{article}

\usepackage[T1]{fontenc}
\usepackage[utf8]{inputenc}
\usepackage[brazilian]{babel}
% margens laterais mínimas para as imagens ocuparem quase toda a página
\usepackage[top=2cm, bottom=2cm, left=1cm, right=1cm]{geometry}
\usepackage{graphicx}
\usepackage{subcaption}
\usepackage{booktabs}
\usepackage{float}
\usepackage{amsmath}
\usepackage{caption}
\usepackage{xcolor}
\usepackage[hidelinks]{hyperref}

\captionsetup{font=small, labelfont=bf}
\captionsetup[sub]{font=footnotesize}

\graphicspath{{figuras_analise_exploratoria/}{figuras_estimacao/}{figuras_fase2/}}

% subfigura padrão: largura quase total, uma por linha, centralizada
\newcommand{\subfig}[2]{%
  \begin{subfigure}[t]{0.96\textwidth}
    \centering
    \includegraphics[width=\linewidth]{#1}
    \caption{#2}
  \end{subfigure}\par\medskip}

% bloco "Como ler": explicação curta e acessível de cada figura
\newcommand{\comoler}[1]{%
  \par\medskip
  \noindent\begingroup\small\itshape
  \textcolor{black!70}{\textbf{Como ler:} #1}%
  \endgroup\par\medskip}

\title{Gestão Hospitalar no SUS}

\begin{document}

% ══════════════════════════════ CAPA ══════════════════════════════
\begin{titlepage}
  \centering
  \vspace*{2.5cm}
  {\LARGE\bfseries Gestão Hospitalar no SUS: o Impacto do Modelo
   Institucional e Organizacional na Prestação de Serviço\par}
  \vspace{1.2cm}
  {\Large Análise Exploratória e Estimação do Painel de Hospitais
   SUS/SP, 2015 a 2025\par}
  \vspace{3.5cm}
  {\large
   João Eduardo Pastori Garcia\\
   {\normalsize Estatístico Responsável}\\[1.4em]
   Priscilla Reinisch Perdicaris\\
   {\normalsize Pesquisadora Principal}\\[1.4em]
   Alberto Tomasi Diniz Tiefensee\\
   {\normalsize Revisor e Responsável pelo Banco de Dados}\par}
  \vfill
  {\large São Paulo\\ Julho de 2026}
\end{titlepage}

\tableofcontents
\clearpage

\section{O painel de dados e como ler este documento}
\label{sec:painel}

Esta parte do documento apresenta a Análise Exploratória dos dados da
pesquisa: antes de qualquer conclusão sobre o efeito do modelo de
gestão, é preciso conhecer em detalhe o material com que trabalhamos.
O que se faz aqui é responder perguntas simples e fundamentais: como se
distribuem os indicadores de desempenho entre os hospitais? Quem são os
hospitais de cada categoria de gestão e eles são comparáveis entre si?
O que aconteceu nos anos da pandemia? O que os dados mostram nos cinco
hospitais que trocaram a administração Direta pela gestão por OSS? E
onde estão os problemas de qualidade de registro que precisam de
cuidado?

\subsection{O painel}

A base é o painel consolidado do projeto, com \textbf{314 hospitais
acompanhados ano a ano de 2015 a 2025, totalizando 3.454 observações}
(uma observação corresponde a um hospital em um ano). O painel é
\textbf{balanceado}: todos os 314 hospitais aparecem nos 11 anos, sem
buracos, o que permite comparar a evolução de cada um deles em bases
iguais. Esse conjunto resulta de um funil documentado de seleção que
partiu de 830 estabelecimentos da base bruta do SIH/SUS e reteve
apenas hospitais com porte mínimo (mais de 50 leitos na mediana dos
anos), produção contínua no período e classificação de complexidade
disponível, excluindo ainda os hospitais de campanha da pandemia.

Os indicadores analisados são os cinco oficiais do estudo, sempre nas
versões que excluem as internações por COVID: \textbf{mortalidade
geral} (proporção de internações que terminam em óbito);
\textbf{mortalidade ajustada}, que exclui os óbitos fetais, maternos e
neonatais; \textbf{tempo médio de permanência} (TMP, em dias);
\textbf{custo por saída} (valor pago por internação encerrada, em
reais); e \textbf{fração de alta complexidade} (proporção das
internações classificadas como de alta complexidade). A eles se somam
as duas \textbf{taxas de ocupação} (internação e UTI), em escala
percentual.

Três regras valem para todo o documento. Primeira: a categoria de
gestão de cada hospital (Direta, OSS, Público Municipal, Filantrópico,
Privado) segue a classificação assistencial do SIH como definição
adotada pela equipe; a categoria \textbf{Privado reúne apenas 3
hospitais} e por isso nunca deve ser lida como retrato de um grupo,
apenas registrada. Segunda: o \textbf{porte} de cada hospital é fixo,
calculado sobre a mediana de leitos de todos os anos com os cortes
oficiais (até 50 leitos, pequeno porte; 51 a 150, médio; 151 a 500,
grande; acima de 500, especial); como o funil já excluiu os pequenos, o
painel tem \textbf{173 hospitais médios, 130 grandes e 11 especiais}, e
o rótulo anual do SIH, que oscila com o cadastro, não é usado.
Terceira: nas análises que envolvem mortalidade, a complexidade do
hospital entra sempre pela versão \textbf{estrutural} do escore
(baseada só em estrutura física e tecnológica), nunca pela versão que
usa a própria mortalidade, para não misturar causa e consequência.

\subsection{Guia rápido de leitura dos gráficos}

Alguns tipos de gráfico e de medida se repetem ao longo do documento;
vale fixá los uma única vez, em linguagem direta.

\textbf{Mediana e média.} A mediana é o valor do hospital do meio da
fila: metade dos hospitais está abaixo dela, metade acima. A média é a
soma dividida pelo total e é sensível a valores extremos: um único
hospital muito caro puxa a média para cima, mas quase não move a
mediana. Como nossos indicadores têm valores extremos legítimos, a
mediana é quase sempre o retrato mais fiel do hospital típico, e por
isso é a medida preferida aqui. Os \textbf{percentis} generalizam a
ideia: o percentil 90 (p90) é o valor abaixo do qual estão 90\% dos
hospitais.

\textbf{Histograma.} É a fotografia da distribuição: o eixo horizontal
traz o valor do indicador e a altura de cada barra mostra quantas
observações caem naquele trecho. Onde as barras são altas, os hospitais
se concentram; caudas compridas para a direita indicam poucos casos com
valores muito altos.

\textbf{Gráfico de violino.} É um histograma de pé e espelhado, um por
categoria: a silhueta engorda onde há mais hospitais. As linhas
internas marcam os quartis e o número impresso é a mediana do grupo.

\textbf{Gráfico QQ.} Compara os dados reais com um molde matemático de
distribuição. Se os pontos seguem a reta, o molde descreve bem os
dados; desvios sistemáticos da reta indicam onde o molde falha. O
número r que acompanha cada QQ resume essa aderência: quanto mais
próximo de 1, melhor.

\textbf{AIC.} É uma nota de qualidade de ajuste usada para comparar
moldes de distribuição entre si: \textbf{quanto menor o AIC, melhor o
molde descreve os dados}. A nota só serve para comparação; o valor
isolado não tem interpretação.

\textbf{Escala logarítmica.} Alguns gráficos usam eixo em escala
logarítmica, em que as marcas avançam por multiplicação (100, 1.000,
10.000) em vez de soma. Isso permite enxergar ao mesmo tempo hospitais
de R\$~500 e de R\$~20.000 por saída sem que os pequenos virem um
borrão; os rótulos dos eixos mostram sempre os valores reais.

\section{Como se distribuem os indicadores}
\label{sec:univariada}

A Tabela~\ref{tab:univariada} resume as distribuições dos sete
indicadores e as Figuras~\ref{fig:univA} a~\ref{fig:univC} trazem os
histogramas. Antes
dos números, um esclarecimento que simplifica todo o restante do
documento: a \textbf{mortalidade ajustada é, na prática, idêntica à
mortalidade geral} neste painel (a exclusão dos óbitos fetais, maternos
e neonatais quase não altera as taxas; a concordância entre as duas é
total, com correlação de 1,000). Por isso ela aparece na tabela, mas as
figuras mostram apenas a mortalidade geral, e tudo o que for dito vale
igualmente para a versão ajustada.

Em palavras simples, o que os dados mostram: o hospital típico do
painel tem \textbf{mortalidade em torno de 5,4\%} (metade dos hospitais
está entre 3,0\% e 7,0\%), \textbf{permanência média de 4,4 dias} e
\textbf{custo de cerca de R\$~1.071 por internação encerrada}. Mas os
extremos importam. Na mortalidade há um segundo grupo de hospitais com
taxas próximas de zero (maternidades e especializados) e \textbf{1,5\%
das observações têm mortalidade exatamente zero}. No tempo de
permanência existe uma \textbf{subpopulação que se acumula num teto de
30,5 dias}: são hospitais de perfil psiquiátrico ou de cuidados
prolongados, em que a permanência longa é característica do serviço (o
teto reflete o ciclo mensal de faturamento da AIH), e não erro de
registro. No custo, a cauda direita é longa: enquanto a mediana é
R\$~1.071, o máximo chega a R\$~27.888, e é por isso que a média
(R\$~1.475) fica 38\% acima da mediana. A fração de alta complexidade é
o indicador mais concentrado: \textbf{34,4\% dos hospitais ano não
realizam nenhuma internação de alta complexidade}. Na ocupação,
\textbf{14,7\% das observações de internação ultrapassam 100\%}
(hospital operando acima da capacidade cadastrada, informação real de
sobrecarga) e a UTI combina \textbf{25,5\% de zeros} (hospitais sem UTI
ativa) com um máximo impossível de 875\%, examinado na
Seção~\ref{sec:qualidade}.

\begin{table}[H]
\centering
\caption{Resumo estatístico dos sete indicadores (3.454 observações de
hospital e ano; mortalidades e alta complexidade em proporção;
ocupações em percentual). DP é o desvio padrão, medida do espalhamento
dos valores; p90 é o valor abaixo do qual estão 90\% das observações;
assimetria positiva indica cauda comprida à direita.}
\label{tab:univariada}
\small
\begin{tabular}{lrrrrrrr}
\toprule
Indicador & Média & DP & Mediana & p90 & Máximo & Assimetria & Zeros \\
\midrule
Mortalidade geral      & 0,0517 & 0,0319 & 0,0536 & 0,0878 & 0,2602 & 0,72 & 1,5\% \\
Mortalidade ajustada   & 0,0517 & 0,0319 & 0,0535 & 0,0878 & 0,2602 & 0,72 & 1,6\% \\
TMP (dias)             & 5,99   & 5,75   & 4,42   & 7,98   & 30,50  & 3,33 & 0,03\% \\
Custo por saída (R\$)  & 1.475,00 & 1.673,59 & 1.071,27 & 2.459,93 & 27.887,60 & 6,59 & 0\% \\
Fração alta complex.   & 0,0731 & 0,1530 & 0,0020 & 0,2215 & 1,0000 & 3,41 & 34,4\% \\
Ocupação internação (\%) & 71,42 & 27,72 & 73,46 & 103,83 & 236,61 & 0,05 & 0\% \\
Ocupação UTI (\%)      & 56,83  & 45,29  & 67,15  & 96,10  & 875,21 & 2,72 & 25,5\% \\
\bottomrule
\end{tabular}
\end{table}

\begin{figure}[H]
\centering
\subfig{fig_ae_01_hist_mort_all.png}{Mortalidade geral}
\subfig{fig_ae_01_hist_tmp.png}{TMP}
\caption{Distribuição da mortalidade geral e do tempo médio de
permanência (314 hospitais, 2015 a 2025). Em cada gráfico, a linha
tracejada vertical é a mediana, a pontilhada é a média e os valores de
pico e máximo estão anotados.}
\label{fig:univA}
\end{figure}

\comoler{cada barra mostra quantas observações caem naquele valor do
indicador. Na mortalidade, note as duas corcovas: a da esquerda, perto
de zero, são as maternidades e especializados; a principal fica perto
de 5,5\%. No TMP, quase todos os hospitais estão entre 3 e 8 dias, mas
há um pequeno grupo isolado junto ao máximo de 30,5 dias, os hospitais
de longa permanência.}

\begin{figure}[H]
\centering
\subfig{fig_ae_01_hist_custo_saida.png}{Custo por saída}
\subfig{fig_ae_01_hist_pct_alta_complex.png}{Fração alta complexidade}
\caption{Distribuição do custo por saída e da fração de alta
complexidade.}
\label{fig:univB}
\end{figure}

\comoler{no custo, a concentração fica perto de R\$~1.000 e a cauda se
estica até quase R\$~28.000, o que explica a média (pontilhada) ficar
bem à direita da mediana (tracejada). Na alta complexidade, a barra
gigante em zero mostra os 34,4\% de hospitais ano sem nenhuma
internação desse tipo.}

\begin{figure}[H]
\centering
\subfig{fig_ae_01_hist_ocupacao_internacao.png}{Ocupação de internação}
\subfig{fig_ae_01_hist_ocupacao_uti.png}{Ocupação de UTI}
\caption{Distribuição das taxas de ocupação de internação e de UTI, em
escala percentual.}
\label{fig:univC}
\end{figure}

\comoler{a linha pontilhada vertical marca 100\%: tudo à direita dela é
operação acima da capacidade cadastrada. Na internação, isso ocorre em
14,7\% das observações. Na UTI, a barra em zero são os hospitais sem
UTI ativa, e a cauda que se arrasta até 875\% é o problema de cadastro
examinado na Seção~\ref{sec:qualidade}.}

\section{Que molde estatístico descreve cada indicador}
\label{sec:distribuicoes}

Cada indicador tem uma forma própria de variar entre hospitais, e há
moldes matemáticos (as chamadas famílias de distribuição) desenhados
para cada tipo de forma. Identificar o molde certo é importante porque
toda análise posterior fica mais confiável quando respeita a natureza
do dado: proporções entre 0 e 1 pedem um molde, valores em reais com
cauda longa pedem outro, contagens pedem um terceiro. Esta seção testa
os moldes contra os dados e documenta as escolhas.

\subsection{Mortalidade: uma proporção com dois grupos de hospitais}

A mortalidade é uma proporção (fica entre 0 e 1), e o molde clássico
para proporções é a distribuição \textbf{Beta}. O ajuste de uma Beta
única descreve bem a escala geral, mas erra o pico de hospitais com
mortalidade quase nula. A Figura~\ref{fig:mistura} mostra a solução e,
ao mesmo tempo, a explicação: somando \textbf{duas Betas} (uma para
cada grupo de hospitais), a curva acompanha o histograma com nitidez, e
a nota de ajuste melhora de forma expressiva (ganho de 1.124 pontos de
AIC). O ponto central para a gestão é que esses dois grupos
\textbf{existem de verdade no painel}: um grupo de \textbf{58 hospitais
de baixo óbito estrutural} (mortalidade mediana abaixo de 2\%,
majoritariamente maternidades, pediátricos e oncológicos) e os
\textbf{256 demais}, hospitais gerais com mortalidade em torno de 6\%.
A aparente anomalia do histograma é, portanto, apenas o retrato de dois
perfis assistenciais diferentes convivendo na rede. Como também existem
observações com mortalidade exatamente zero (1,5\%), o molde completo
para a mortalidade é a \textbf{Beta com um tratamento especial para os
zeros}.

\begin{figure}[H]
\centering
\subfig{fig_ae_12_mistura_mort_all.png}{Um molde só vs a soma de dois}
\subfig{fig_ae_12_decomposicao_mort_all.png}{Os dois grupos reais de
hospitais}
\caption{Mortalidade geral: à esquerda, a curva cinza tracejada é o
molde único (que erra o pico perto de zero) e a curva escura é a soma
de dois moldes, que acompanha o histograma; à direita, a mesma soma
reconstruída a partir dos dois grupos reais de hospitais do painel.}
\label{fig:mistura}
\end{figure}

\comoler{no painel da esquerda, compare as duas curvas com as barras: a
tracejada passa por cima do vale e erra os dois picos; a escura desenha
os dois picos. No painel da direita, as curvas tracejadas coloridas são
os dois grupos de hospitais (baixo óbito à esquerda, demais à direita)
e a curva preta é a soma dos dois, praticamente idêntica ao histograma:
a forma estranha da mortalidade é só a mistura de dois perfis de
hospital.}

\subsection{Alta complexidade: o peso dos zeros}

A fração de alta complexidade também é uma proporção, mas com uma
particularidade dominante: \textbf{um terço das observações vale
exatamente zero} (hospitais que não realizam alta complexidade) e uma
pequena parcela vale exatamente um (hospitais que só realizam alta
complexidade). O molde adequado trata três situações em separado: a
chance de o hospital estar no zero, a chance de estar no um, e a
distribuição de quem está no meio; esse molde combinado é conhecido
pela sigla \textbf{ZOIB}. A Figura~\ref{fig:betapct} mostra o ajuste do
miolo e a Figura~\ref{fig:zeros} mostra que \textbf{os zeros não são
aleatórios}: a proporção de hospitais sem alta complexidade cai de
82,6\% na faixa 2 de complexidade para zero na faixa 6, e de 41,1\% nos
Filantrópicos para 10,3\% nas OSS. Ou seja, saber a complexidade, o
porte e a categoria do hospital já diz muito sobre a chance de ele
estar no grupo do zero.

\begin{figure}[H]
\centering
\includegraphics[width=.96\textwidth]{fig_ae_02_beta_pct_alta_complex.png}
\caption{Fração de alta complexidade entre os hospitais que realizam
alguma alta complexidade: histograma e molde Beta ajustado. O eixo
vertical está em escala logarítmica para a cauda inteira ficar
visível.}
\label{fig:betapct}
\end{figure}

\comoler{a grande maioria dos hospitais que fazem alta complexidade faz
pouca (a concentração à esquerda; metade fica abaixo de 3,7\%), e uma
minoria é altamente especializada (as barras à direita, perto de
100\%). A curva escura é o molde Beta, que acompanha essa queda
contínua.}

\begin{figure}[H]
\centering
\subfig{fig_ae_15_zeros_pct_alta_categoria.png}{Por categoria de
gestão}
\subfig{fig_ae_15_zeros_pct_alta_faixa.png}{Por faixa de complexidade}
\caption{Percentual de observações com alta complexidade exatamente
zero, por grupo de hospitais.}
\label{fig:zeros}
\end{figure}

\comoler{cada barra é a proporção de hospitais ano daquele grupo que
não realizou nenhuma internação de alta complexidade. A queda das
barras da esquerda para a direita no painel (b) mostra que o zero é
questão de perfil e estrutura: quase todos os hospitais de faixa 2
estão no zero, e nenhum de faixa 6 está.}

\subsection{Faturamento e tempo de permanência: valores positivos com cauda
longa}

Faturamento por saída e TMP são valores positivos com cauda à direita. Cinco
moldes clássicos foram postos à prova
(Tabela~\ref{tab:familias})\footnote{Guia dos tamanhos de amostra que
aparecem no documento: \textbf{3.454} é o painel completo;
\textbf{3.453} exclui a única observação com TMP igual a zero (Hospital
Infante D.\ Henrique, 2021), incompatível com a escala logarítmica;
\textbf{3.255} exclui também as 198 observações dos 18 hospitais de
longa permanência; \textbf{3.452} exclui as duas observações de
denominador frágil da Seção~\ref{sec:qualidade}, regra usada nos
modelos da fase de estimação; e \textbf{2.571} são as observações com
UTI ativa depois de retiradas essas duas frágeis. Cada conta foi
conferida no painel.},
comparados pela nota AIC e pelo gráfico QQ. O resultado é consistente:
o molde \textbf{log logístico (também chamado de Fisk)} vence nas três
amostras, com a \textbf{LogNormal} logo atrás, e a Gama, molde usado
por tradição em análises de custo, fica sempre entre as piores. No TMP
há um detalhe decisivo: enquanto os hospitais de longa permanência
estiverem na amostra, nenhum molde acompanha o teto de 30,5 dias;
\textbf{retirando os 18 hospitais de longa permanência, o ajuste do
molde log logístico fica praticamente perfeito} (r de 0,9975 na
Figura~\ref{fig:qqtmpsemlp}). A leitura prática: faturamento e TMP devem ser
analisados na escala logarítmica, e os hospitais de longa permanência
merecem tratamento em separado nas análises de TMP. Os gráficos QQ
completos estão no Apêndice~\ref{ap:diagnosticos}: no faturamento, o molde
log logístico acompanha os dados por quase toda a extensão (r de
0,9754, Figura~\ref{fig:qqcusto}); no TMP do painel completo, nenhum
molde reproduz o teto de 30,5 dias (Figura~\ref{fig:qqtmp}).

\begin{table}[H]
\centering
\caption{Comparação dos cinco moldes de distribuição para faturamento e TMP
pela nota AIC (Seção~\ref{sec:painel}) e pelo r do QQ
(Apêndice~\ref{ap:diagnosticos}).}
\label{tab:familias}
\small
\begin{tabular}{llrrr}
\toprule
Amostra & Molde & $n$ & AIC & r do QQ \\
\midrule
Faturamento por saída & Log logística (Fisk) & 3.454 & 55.483 & 0,9754 \\
                & LogNormal            & 3.454 & 55.547 & 0,9162 \\
                & Gaussiana inversa    & 3.454 & 55.608 & 0,9018 \\
                & Gama                 & 3.454 & 56.499 & 0,8411 \\
                & Weibull              & 3.454 & 56.995 & 0,8415 \\
\midrule
TMP, painel completo & Log logística (Fisk) & 3.453 & 15.990 & 0,8597 \\
                     & LogNormal            & 3.453 & 16.650 & 0,8696 \\
                     & Gaussiana inversa    & 3.453 & 16.767 & 0,8736 \\
                     & Gama                 & 3.453 & 17.990 & 0,8379 \\
                     & Weibull              & 3.453 & 18.701 & 0,8399 \\
\midrule
TMP sem longa permanência & Log logística (Fisk) & 3.255 & 12.147 & 0,9975 \\
                          & LogNormal            & 3.255 & 12.226 & 0,9810 \\
                          & Gaussiana inversa    & 3.255 & 12.260 & 0,9787 \\
                          & Gama                 & 3.255 & 12.417 & 0,9635 \\
                          & Weibull              & 3.255 & 13.204 & 0,9403 \\
\bottomrule
\end{tabular}
\end{table}

\begin{figure}[H]
\centering
\subfig{fig_ae_03_hist_log_custo_saida.png}{Faturamento por saída}
\subfig{fig_ae_03_hist_log_tmp.png}{TMP}
\caption{Distribuições de faturamento e TMP com o eixo horizontal em escala
logarítmica (os rótulos mostram os valores reais).}
\label{fig:histlog}
\end{figure}

\comoler{na escala logarítmica, o faturamento (a) vira um morro quase
simétrico, sinal clássico de que a análise deve ser feita nessa escala.
No TMP (b), além do morro principal entre 3 e 8 dias, aparece um
segundo morrinho isolado à direita, junto de 30 dias: é a subpopulação
de longa permanência.}

\begin{figure}[H]
\centering
\includegraphics[width=.96\textwidth]{fig_ae_03_qq_fisk_tmp_sem_longa_perm.png}
\caption{Gráfico QQ do TMP sem os 18 hospitais de longa permanência: os
pontos ficam praticamente sobre a reta (r de 0,9975).}
\label{fig:qqtmpsemlp}
\end{figure}

\comoler{este é o gráfico QQ do TMP refeito sem os 18 hospitais de
longa permanência (a versão com o painel completo está na
Figura~\ref{fig:qqtmp}, no apêndice): o achatamento desaparece e o
molde descreve os
dados quase perfeitamente. É a demonstração visual de que o problema
não era o molde, era a mistura de dois perfis de hospital.}

\subsection{Produção: contagens muito mais variáveis que o usual}

A produção anual (número de saídas) é uma contagem. O molde mais
simples para contagens, a Poisson, supõe que a variação entre hospitais
seja comparável à média; aqui \textbf{a variação é milhares de vezes
maior que a média} (razão de 6.381 no agregado e de 93 mesmo comparando
cada hospital só consigo mesmo ao longo dos anos). O molde adequado é a
\textbf{Binomial Negativa}, que aceita essa variação extra.

\subsection{Os moldes melhoram quando se compara o comparável}

Uma pergunta natural: se os QQ do painel completo não são perfeitos, o
molde está errado? As Figuras~\ref{fig:qqcondA} a~\ref{fig:qqcondC},
reunidas no Apêndice~\ref{ap:diagnosticos},
respondem que não: o desajuste vem de misturar 314 hospitais muito
diferentes num único gráfico. Quando se compara cada hospital consigo mesmo (controlando o
efeito do hospital e do ano), os desvios praticamente desaparecem: no
faturamento, a assimetria restante é quase nula; na mortalidade, a forte
assimetria da distribuição bruta na escala \textbf{logito} (a escala
interna em que o modelo trabalha proporções), de 1,72 negativa, cai
para 0,11 depois do controle; na escala original da
Tabela~\ref{tab:univariada} a assimetria é outra (0,72 positiva), e os
dois números não se contradizem. Em resumo, \textbf{os moldes escolhidos descrevem bem os
dados quando se compara o comparável}.

\section{O perfil de cada categoria de gestão}
\label{sec:categorias}

O painel contém \textbf{187 hospitais Filantrópicos, 58 Públicos
Municipais, 39 sob gestão de OSS, 32 da administração Direta e 3
Privados}; as contagens de OSS e Direta variam levemente no tempo por
causa dos cinco hospitais que trocaram de gestão. A
Tabela~\ref{tab:categorias} traz as medianas de cada grupo e as
Figuras~\ref{fig:violA} a~\ref{fig:violC} mostram a distribuição
completa.

O retrato operacional é nítido e fala direto à gestão: \textbf{os
hospitais OSS operam com ocupação mediana de 91,6\% e os Públicos
Municipais de 87,3\%, contra 68,2\% dos Filantrópicos e apenas 62,8\%
da administração Direta}. O tempo de permanência conta história
parecida: \textbf{6,07 dias na Direta contra 4,85 nas OSS}, 4,77 nos
Municipais e 4,11 nos Filantrópicos. Já na mortalidade a ordem se
inverte: Direta (4,75\%) e OSS (4,93\%) ficam abaixo de Filantrópicos
(5,45\%) e Municipais (5,64\%). É essencial não tirar conclusões de
causa dessas comparações brutas: as categorias atendem perfis
diferentes de complexidade e porte, e essas diferenças precisam ser
descontadas antes de qualquer julgamento. O custo mediano por saída
varia pouco entre as categorias públicas e filantrópicas (de R\$~979
nos Municipais a R\$~1.193 na Direta). A categoria \textbf{Privado é um
mundo à parte}: custo mediano de R\$~12.798 e 98,4\% de alta
complexidade, perfil de centros especializados; com só 3 hospitais,
serve de registro, não de comparação.

\begin{table}[H]
\centering
\caption{Medianas por categoria de gestão, período completo
(mortalidade em proporção; ocupações em percentual).}
\label{tab:categorias}
\small
\begin{tabular}{lrrrrrr}
\toprule
Categoria & Mortalidade & TMP & Custo (R\$) & Alta complex. &
Ocup.\ int.\ (\%) & Ocup.\ UTI (\%) \\
\midrule
Direta            & 0,0475 & 6,07 & 1.192,67  & 0,0088 & 62,8 & 57,3 \\
OSS               & 0,0493 & 4,85 & 1.147,40  & 0,0242 & 91,6 & 73,4 \\
Público Municipal & 0,0564 & 4,77 & 978,96    & 0,0010 & 87,3 & 69,7 \\
Filantrópico      & 0,0545 & 4,11 & 1.086,13  & 0,0011 & 68,2 & 66,8 \\
Privado           & 0,0256 & 4,24 & 12.798,36 & 0,9836 & 40,5 & 3,5  \\
\bottomrule
\end{tabular}
\end{table}

\begin{figure}[H]
\centering
\subfig{fig_ae_04_violino_mort_all.png}{Mortalidade geral}
\subfig{fig_ae_04_violino_tmp.png}{TMP (escala log)}
\caption{Distribuição da mortalidade e do TMP por categoria de gestão.
Cada violino é a silhueta da distribuição do grupo; o número impresso é
a mediana.}
\label{fig:violA}
\end{figure}

\comoler{a silhueta engorda onde há mais hospitais; as linhas internas
marcam os quartis e o número é a mediana. Na mortalidade, os violinos
das quatro categorias públicas e filantrópicas ocupam alturas
parecidas; no TMP, o violino da Direta está visivelmente deslocado para
cima.}

\begin{figure}[H]
\centering
\subfig{fig_ae_04_violino_custo_saida.png}{Custo por saída (escala
log)}
\subfig{fig_ae_04_violino_pct_alta_complex.png}{Fração alta
complexidade}
\caption{Distribuição do custo por saída e da fração de alta
complexidade por categoria de gestão.}
\label{fig:violB}
\end{figure}

\comoler{no custo, os quatro grupos públicos e filantrópicos são
parecidos e o Privado flutua sozinho lá em cima, uma ordem de grandeza
acima. Nos Filantrópicos, os violinos mais compridos refletem a enorme
diversidade interna do grupo, das pequenas Santas Casas aos grandes
hospitais de referência.}

\begin{figure}[H]
\centering
\subfig{fig_ae_04_violino_ocupacao_internacao.png}{Ocupação de
internação}
\subfig{fig_ae_04_violino_ocupacao_uti.png}{Ocupação de UTI}
\caption{Distribuição das taxas de ocupação por categoria de gestão.}
\label{fig:violC}
\end{figure}

\comoler{na ocupação de internação, veja como os violinos de OSS e
Municipal estão inteiros mais altos que os de Direta e Filantrópico: é
a diferença de regime operacional em uma imagem. Na UTI, os corpos dos
violinos são mais parecidos, mas as caudas de zeros (hospitais sem UTI)
puxam as silhuetas para baixo.}

Uma pergunta indispensável antes de comparar categorias:
\textbf{elas têm hospitais parecidos o bastante para a comparação ser
justa?} A Figura~\ref{fig:composicao} responde contando os hospitais de
cada categoria em cada faixa de complexidade e porte. A sobreposição é
boa nas faixas 3 e 4, onde estão 24 dos 31 hospitais da Direta e 31 dos
35 das OSS; mas \textbf{a Direta não tem nenhum hospital na faixa 5 e
as OSS não têm nenhum nas faixas 2 e 6}; o Privado só existe nas faixas
2 e 3; e a faixa 2 é quase toda Filantrópica (37 dos 49 hospitais). No
porte, as OSS não têm hospital especial. Em consequência,
\textbf{comparações entre categorias são seguras nas faixas 3 e 4 e nos
portes médio e grande}, onde todos os grupos públicos e filantrópicos
estão presentes; fora dessa região faltam hospitais para comparar.

\begin{figure}[H]
\centering
\subfig{fig_ae_13_composicao_categoria_faixa.png}{Categoria por faixa
de complexidade}
\subfig{fig_ae_13_composicao_categoria_porte.png}{Categoria por porte}
\caption{Quantos hospitais cada categoria tem em cada faixa de
complexidade e em cada porte (contagem de CNES; azul mais escuro
indica mais hospitais).}
\label{fig:composicao}
\end{figure}

\comoler{cada célula informa o número de hospitais daquela categoria
naquela faixa ou porte. As células com zero são o alerta: onde há zero,
não existe hospital daquele perfil naquela categoria, e portanto não há
como comparar. A mancha escura na linha dos Filantrópicos, faixas 2 e
3, mostra onde está o grosso da rede.}

\section{Complexidade e porte: os grandes divisores da rede}
\label{sec:complexidadeporte}

Antes de qualquer comparação entre modelos de gestão, é preciso
reconhecer o quanto a complexidade e o porte, sozinhos, separam os
hospitais. As Figuras~\ref{fig:faixaA} a~\ref{fig:faixaC} mostram a
mediana de cada indicador ao longo das faixas Barcelona de complexidade
(da faixa 2, menor, à 6, maior), e as Figuras~\ref{fig:porteA}
a~\ref{fig:porteC}, ao longo dos portes.

O padrão é de escada: o custo mediano por saída sobe de \textbf{R\$~495
na faixa 2 para R\$~2.772 na faixa 6} (mais de cinco vezes), o tempo de
permanência vai de 3,27 para 6,30 dias, a fração de alta complexidade
sai de zero e chega a 22,5\%, e a ocupação de internação sobe de 37,5\%
para a casa dos 90\%. A mortalidade cresce com a complexidade até a
faixa 5 (de 2,88\% para 6,40\%) e \textbf{recua na faixa 6 (5,43\%)},
um sinal de que os hospitais de ponta convertem estrutura em resultado.
Vale notar que metade dos 49 hospitais da faixa 2 não opera UTI. No
porte, a mesma escada: custo mediano de \textbf{R\$~834 nos médios,
R\$~1.470 nos grandes e R\$~2.754 nos especiais}.

A conclusão para todo o restante do documento: \textbf{complexidade e
porte precisam ser sempre levados em conta ao comparar categorias de
gestão}, sob pena de atribuir ao modelo de gestão diferenças que são,
na verdade, de perfil assistencial. E, pela salvaguarda definida na
Seção~\ref{sec:painel}, a complexidade entra nas análises de
mortalidade apenas pela versão estrutural do escore.

\begin{figure}[H]
\centering
\subfig{fig_ae_05_faixa_mort_all.png}{Mortalidade geral}
\subfig{fig_ae_05_faixa_tmp.png}{TMP}
\caption{Mediana da mortalidade e do TMP por faixa Barcelona de
complexidade (2 a 6), com o valor sobre cada barra.}
\label{fig:faixaA}
\end{figure}

\comoler{cada barra é a mediana do indicador naquela faixa de
complexidade, com o valor impresso em cima. Na mortalidade, repare no
degrau que sobe até a faixa 5 e desce na 6; no TMP, a escada é sempre
ascendente.}

\begin{figure}[H]
\centering
\subfig{fig_ae_05_faixa_custo_saida.png}{Custo por saída}
\subfig{fig_ae_05_faixa_pct_alta_complex.png}{Fração alta complexidade}
\caption{Mediana do custo por saída e da fração de alta complexidade
por faixa Barcelona.}
\label{fig:faixaB}
\end{figure}

\comoler{a escada do custo é a mais íngreme do documento: da faixa 2
para a 6, a mediana mais que quintuplica. A alta complexidade parte do
zero absoluto nas faixas baixas.}

\begin{figure}[H]
\centering
\subfig{fig_ae_05_faixa_ocupacao_internacao.png}{Ocupação de
internação}
\subfig{fig_ae_05_faixa_ocupacao_uti.png}{Ocupação de UTI}
\caption{Mediana das taxas de ocupação por faixa Barcelona.}
\label{fig:faixaC}
\end{figure}

\comoler{quanto mais complexo o hospital, mais cheio ele opera. Na
UTI, a barra da faixa 2 vale zero porque metade daqueles hospitais não
tem UTI.}

\begin{figure}[H]
\centering
\subfig{fig_ae_06_porte_mort_all.png}{Mortalidade geral}
\subfig{fig_ae_06_porte_tmp.png}{TMP}
\caption{Mediana da mortalidade e do TMP por porte (classificação fixa
pela mediana de leitos: 173 hospitais médios, 130 grandes e 11
especiais).}
\label{fig:porteA}
\end{figure}

\comoler{mesma leitura das figuras por faixa, agora por porte:
hospitais maiores seguram o paciente por mais tempo e têm mortalidade
um pouco maior, coerente com casos mais graves.}

\begin{figure}[H]
\centering
\subfig{fig_ae_06_porte_custo_saida.png}{Custo por saída}
\subfig{fig_ae_06_porte_pct_alta_complex.png}{Fração alta complexidade}
\caption{Mediana do custo por saída e da fração de alta complexidade
por porte.}
\label{fig:porteB}
\end{figure}

\comoler{o custo mediano dos hospitais especiais é mais de três vezes o
dos médios, e a alta complexidade se concentra quase toda nos portes
grande e especial.}

\begin{figure}[H]
\centering
\subfig{fig_ae_06_porte_ocupacao_internacao.png}{Ocupação de
internação}
\subfig{fig_ae_06_porte_ocupacao_uti.png}{Ocupação de UTI}
\caption{Mediana das taxas de ocupação por porte.}
\label{fig:porteC}
\end{figure}

\comoler{hospitais maiores operam mais cheios, tanto na internação
quanto na UTI: mais um retrato de que porte e pressão assistencial
andam juntos.}

\section{A evolução no tempo e a marca da pandemia}
\label{sec:tendencia}

As Figuras~\ref{fig:tendA} a~\ref{fig:tendC} acompanham a mediana anual
de cada indicador de 2015 a 2025, com a faixa cinza destacando o biênio
2020 e 2021. A pandemia deixou marca em quase tudo: a mortalidade (já sem as
internações COVID) sobe de cerca de 5,1\% para \textbf{5,63\% em 2020 e
5,94\% em 2021} e só retorna ao patamar histórico em 2025, sinal de que
o biênio afetou também os pacientes não COVID. A ocupação de internação
despenca para \textbf{65,9\% em 2020} (leitos esvaziados por
cancelamentos e medo) enquanto a de UTI dispara para \textbf{82,9\% em
2021}, seu recorde na série. O tempo de permanência atravessa tudo
estável, em torno de 4,4 dias.

O faturamento pede uma leitura em duas camadas. Em valores correntes
(nominais), a mediana sobe de R\$~908 em 2015 para R\$~1.317 em 2025.
Mas parte dessa alta é só inflação. Corrigindo tudo para o poder de
compra de 2025 com o IPCA (Figura~\ref{fig:custoreal} e
Tabela~\ref{tab:custoreal}), a história muda de sinal: \textbf{em
valores de 2025, o faturamento mediano caiu de R\$~1.496 para R\$~1.317, uma
queda real de cerca de 12\% na década}. Toda a alta aparente é
inflação. Qualquer avaliação da evolução do faturamento deve, portanto, ser
feita em valores reais ou com o nível de preços de cada ano devidamente
considerado.

Por categoria (Figuras~\ref{fig:tendcatA} a~\ref{fig:tendcatC}), o
choque pandêmico é
generalizado, mas a queda de ocupação de 2020 é mais funda na Direta e
nos Filantrópicos; a fração de alta complexidade das OSS cresce de
forma consistente na segunda metade do período; e, no faturamento, a Direta
parte de nível parecido com as demais categorias e passa a liderar o
faturamento mediano a partir de 2020.

\begin{table}[H]
\centering
\caption{Faturamento por saída mediano por ano: em valores correntes e em
valores de 2025 (corrigidos pelo IPCA, variação de dezembro a dezembro,
IBGE). O fator indica por quanto o valor do ano é multiplicado para
chegar a preços de 2025.}
\label{tab:custoreal}
\small
\begin{tabular}{lrrr}
\toprule
Ano & Nominal (R\$) & Valores de 2025 (R\$) & Fator IPCA \\
\midrule
2015 & 907,60   & 1.495,61 & 1,6479 \\
2016 & 907,18   & 1.406,45 & 1,5504 \\
2017 & 921,69   & 1.388,00 & 1,5059 \\
2018 & 931,01   & 1.351,36 & 1,4515 \\
2019 & 957,47   & 1.332,35 & 1,3915 \\
2020 & 1.060,25 & 1.411,56 & 1,3313 \\
2021 & 1.093,68 & 1.322,98 & 1,2097 \\
2022 & 1.220,00 & 1.395,02 & 1,1435 \\
2023 & 1.236,87 & 1.351,85 & 1,0930 \\
2024 & 1.300,81 & 1.356,22 & 1,0426 \\
2025 & 1.316,94 & 1.316,94 & 1,0000 \\
\bottomrule
\end{tabular}
\end{table}

\begin{figure}[H]
\centering
\subfig{fig_ae_07_tendencia_mort_all.png}{Mortalidade geral}
\subfig{fig_ae_07_tendencia_tmp.png}{TMP}
\caption{Mediana anual da mortalidade e do TMP, com a metade central
dos hospitais na área sombreada em azul; a faixa cinza destaca 2020 e
2021 e os pontos de máximo, mínimo e último ano estão anotados.}
\label{fig:tendA}
\end{figure}

\comoler{a linha é o hospital mediano em cada ano; a área azul em volta
contém a metade central dos hospitais (do quartil inferior ao
superior). Dentro da faixa cinza da pandemia, note a corcova da
mortalidade; o TMP, em contraste, atravessa a década quase imóvel.}

\begin{figure}[H]
\centering
\subfig{fig_ae_07_tendencia_custo_saida.png}{Faturamento por saída (valores
correntes)}
\subfig{fig_ae_07_tendencia_pct_alta_complex.png}{Fração alta
complexidade}
\caption{Mediana anual do faturamento por saída (valores correntes) e da
fração de alta complexidade.}
\label{fig:tendB}
\end{figure}

\comoler{o faturamento em valores correntes sobe a década inteira, mas essa
leitura é enganosa; a correção pela inflação está na
Figura~\ref{fig:custoreal}. A alta complexidade mediana cresce nos anos
finais da série.}

\begin{figure}[H]
\centering
\subfig{fig_ae_07_tendencia_ocupacao_internacao.png}{Ocupação de
internação}
\subfig{fig_ae_07_tendencia_ocupacao_uti.png}{Ocupação de UTI}
\caption{Mediana anual das taxas de ocupação.}
\label{fig:tendC}
\end{figure}

\comoler{dentro da faixa cinza da pandemia, os dois movimentos em
sentidos opostos: a internação esvazia em 2020 (mergulho para 65,9\%) e
a UTI lota em 2021 (pico de 82,9\%, recorde da série).}

\begin{figure}[H]
\centering
\includegraphics[width=.96\textwidth]{fig_ae_07_tendencia_custo_saida_real.png}
\caption{Faturamento por saída mediano: a linha clara tracejada é o valor
corrente e a linha escura é o valor corrigido pelo IPCA para o poder de
compra de 2025.}
\label{fig:custoreal}
\end{figure}

\comoler{as duas linhas contam histórias opostas. A clara (valores
correntes) sobe o tempo todo; a escura (valores de 2025) mostra o que o
dinheiro de fato compra e cai ao longo da década: de R\$~1.496 para
R\$~1.317 por internação. Quando as duas linhas se encontram, em 2025,
é porque ali o valor corrente e o corrigido coincidem por definição.}

\begin{figure}[H]
\centering
\subfig{fig_ae_08_categoria_mort_all.png}{Mortalidade geral}
\subfig{fig_ae_08_categoria_tmp.png}{TMP}
\caption{Mediana anual da mortalidade e do TMP por categoria de gestão;
cada categoria tem cor e marcador próprios, e a linha pontilhada cinza
do Privado (3 hospitais) é apenas registro.}
\label{fig:tendcatA}
\end{figure}

\comoler{cada linha é o hospital mediano de uma categoria em cada ano.
Na mortalidade, todas as categorias fazem a corcova da pandemia. No
TMP, a linha da Direta corre acima das demais durante a década
inteira.}

\begin{figure}[H]
\centering
\subfig{fig_ae_08_categoria_custo_saida.png}{Faturamento por saída (Privado
omitido por escala)}
\subfig{fig_ae_08_categoria_pct_alta_complex.png}{Fração alta
complexidade (Privado omitido por escala)}
\caption{Mediana anual do faturamento por saída e da fração de alta
complexidade por categoria. O Privado foi omitido destes dois painéis
porque sua escala (faturamento dez vezes maior) esconderia as demais
categorias; seus valores estão na Tabela~\ref{tab:categorias}.}
\label{fig:tendcatB}
\end{figure}

\comoler{no faturamento, as quatro categorias sobem juntas até 2019 e se
abrem a partir de 2020, com a Direta na frente; lembre que são valores
correntes (a correção pela inflação está na
Figura~\ref{fig:custoreal}). Na alta complexidade, a linha das OSS se
descola das demais na segunda metade do período.}

\begin{figure}[H]
\centering
\subfig{fig_ae_08_categoria_ocupacao_internacao.png}{Ocupação de
internação}
\subfig{fig_ae_08_categoria_ocupacao_uti.png}{Ocupação de UTI}
\caption{Mediana anual das taxas de ocupação por categoria de gestão.}
\label{fig:tendcatC}
\end{figure}

\comoler{na ocupação de internação, as duas linhas de cima (OSS e
Municipal) e as duas de baixo (Filantrópico e Direta) mantêm a
separação de quase 25 pontos percentuais durante toda a década, com o
mergulho de 2020 visível em todas. Na UTI, o pico coletivo de 2021 é a
marca da pandemia.}

\section{Os cinco hospitais que mudaram de gestão}
\label{sec:conversores}

Cinco hospitais do painel trocaram a administração Direta pela gestão
por OSS dentro do período de estudo, e eles são o material mais valioso
da pesquisa: permitem observar o mesmo hospital antes e depois da
mudança. São eles o \textbf{Conjunto Hospitalar de Sorocaba (CNES
2081695)}, convertido em 2019; o \textbf{Centro de Referência da Saúde
da Mulher, Pérola Byington (CNES 2078287)}, convertido em 2023; e três
convertidos em 2025: o \textbf{Hospital Kátia de Souza Rodrigues,
Taipas (CNES 2082225)}, o \textbf{Hospital Geral de Vila Penteado (CNES
2091755)} e o \textbf{Hospital Estadual de Presidente Prudente (CNES
2750511)}. As Figuras~\ref{fig:conv2081695}
a~\ref{fig:conv2750511b} mostram, para cada um, a trajetória de quatro
indicadores contra a referência típica das categorias Direta e OSS.

No \textbf{Conjunto Hospitalar de Sorocaba}, único caso com sete anos
de observação após a mudança, a mortalidade vinha subindo antes da
conversão, atingiu 9,9\% em 2018 e caiu de forma sustentada depois,
chegando a 7,4\% em 2022 antes de um repique recente. Comparando as
medianas antes e depois: mortalidade \textbf{14,8\% menor} (de 9,34\%
para 7,95\%), permanência \textbf{14,0\% menor} (de 8,29 para 7,13
dias) e custo por saída 19,3\% maior, alta parecida com a do setor como
um todo no mesmo intervalo. O hospital segue com mortalidade bem acima
da OSS típica, coerente com seu papel de grande complexo regional que
concentra casos graves.

No \textbf{Pérola Byington}, a mudança de 2023 coincide com melhora
mais expressiva: mortalidade \textbf{39,5\% menor} (de 7,55\% para
4,57\%, convergindo para o padrão das OSS), permanência 17,5\% menor e
custo \textbf{11,4\% menor num período em que o custo do setor subia}.

Duas cautelas de leitura. Em Sorocaba, o pico de mortalidade
imediatamente antes da conversão levanta a possibilidade de a queda
posterior ser, em parte, retorno natural ao patamar anterior. No
Pérola, a melhora começa visivelmente em 2021 e 2022, antes da mudança
formal de gestão, o que sugere que outros fatores também atuavam.
Quanto aos três convertidos em 2025, há \textbf{apenas um ano de
observação após a mudança}, insuficiente para qualquer leitura de
tendência; suas figuras ficam como registro do ponto de partida.

\begin{figure}[H]
\centering
\subfig{fig_ae_09_conversor_2081695_mort_all.png}{Mortalidade geral}
\subfig{fig_ae_09_conversor_2081695_tmp.png}{TMP}
\caption{Conjunto Hospitalar de Sorocaba (CNES 2081695): conversão para
OSS em 2019. Mortalidade geral e TMP.}
\label{fig:conv2081695}
\end{figure}

\begin{figure}[H]
\centering
\subfig{fig_ae_09_conversor_2081695_custo_saida.png}{Custo por saída}
\subfig{fig_ae_09_conversor_2081695_pct_alta_complex.png}{Fração alta
complexidade}
\caption{Conjunto Hospitalar de Sorocaba (CNES 2081695): conversão para
OSS em 2019. Custo por saída e fração de alta complexidade.}
\label{fig:conv2081695b}
\end{figure}

\comoler{em cada painel, a linha preta é o hospital; as linhas
tracejadas coloridas são o hospital mediano da Direta e o da OSS, como
referência; e a linha vertical tracejada marca a entrada na gestão OSS.
Veja na mortalidade a subida até 2018, a queda sustentada depois da
linha vertical e a distância que ainda separa o hospital da OSS
típica.}

\begin{figure}[H]
\centering
\subfig{fig_ae_09_conversor_2078287_mort_all.png}{Mortalidade geral}
\subfig{fig_ae_09_conversor_2078287_tmp.png}{TMP}
\caption{Centro de Referência da Saúde da Mulher, Pérola Byington (CNES
2078287): conversão para OSS em 2023. Mortalidade geral e TMP.}
\label{fig:conv2078287}
\end{figure}

\begin{figure}[H]
\centering
\subfig{fig_ae_09_conversor_2078287_custo_saida.png}{Custo por saída}
\subfig{fig_ae_09_conversor_2078287_pct_alta_complex.png}{Fração alta
complexidade}
\caption{Centro de Referência da Saúde da Mulher, Pérola Byington (CNES
2078287): conversão para OSS em 2023. Custo por saída e fração de alta complexidade.}
\label{fig:conv2078287b}
\end{figure}

\comoler{mesma leitura da figura anterior. Na mortalidade, a queda
começa em 2021, antes da linha vertical, e depois da conversão o
hospital passa a andar colado na referência das OSS; no custo, a linha
preta desce enquanto as referências sobem.}

\begin{figure}[H]
\centering
\subfig{fig_ae_09_conversor_2082225_mort_all.png}{Mortalidade geral}
\subfig{fig_ae_09_conversor_2082225_tmp.png}{TMP}
\caption{Hospital Kátia de Souza Rodrigues, Taipas (CNES 2082225):
conversão para OSS em 2025. Mortalidade geral e TMP.}
\label{fig:conv2082225}
\end{figure}

\begin{figure}[H]
\centering
\subfig{fig_ae_09_conversor_2082225_custo_saida.png}{Custo por saída}
\subfig{fig_ae_09_conversor_2082225_pct_alta_complex.png}{Fração alta
complexidade}
\caption{Hospital Kátia de Souza Rodrigues, Taipas (CNES 2082225):
conversão para OSS em 2025. Custo por saída e fração de alta complexidade.}
\label{fig:conv2082225b}
\end{figure}

\comoler{a linha vertical está no último ano do painel: tudo o que se
vê antes dela é o hospital ainda sob administração Direta. A figura
serve de ponto de partida para acompanhamento futuro, não de avaliação
da mudança.}

\begin{figure}[H]
\centering
\subfig{fig_ae_09_conversor_2091755_mort_all.png}{Mortalidade geral}
\subfig{fig_ae_09_conversor_2091755_tmp.png}{TMP}
\caption{Hospital Geral de Vila Penteado (CNES 2091755): conversão para
OSS em 2025. Mortalidade geral e TMP.}
\label{fig:conv2091755}
\end{figure}

\begin{figure}[H]
\centering
\subfig{fig_ae_09_conversor_2091755_custo_saida.png}{Custo por saída}
\subfig{fig_ae_09_conversor_2091755_pct_alta_complex.png}{Fração alta
complexidade}
\caption{Hospital Geral de Vila Penteado (CNES 2091755): conversão para
OSS em 2025. Custo por saída e fração de alta complexidade.}
\label{fig:conv2091755b}
\end{figure}

\comoler{mesma situação do hospital anterior: um único ano após a
mudança. Note, ainda assim, a subida de mortalidade e de custo ao longo
da década sob gestão Direta, o contexto em que a conversão foi
decidida.}

\begin{figure}[H]
\centering
\subfig{fig_ae_09_conversor_2750511_mort_all.png}{Mortalidade geral}
\subfig{fig_ae_09_conversor_2750511_tmp.png}{TMP}
\caption{Hospital Estadual de Presidente Prudente (CNES 2750511):
conversão para OSS em 2025. Mortalidade geral e TMP.}
\label{fig:conv2750511}
\end{figure}

\begin{figure}[H]
\centering
\subfig{fig_ae_09_conversor_2750511_custo_saida.png}{Custo por saída}
\subfig{fig_ae_09_conversor_2750511_pct_alta_complex.png}{Fração alta
complexidade}
\caption{Hospital Estadual de Presidente Prudente (CNES 2750511):
conversão para OSS em 2025. Custo por saída e fração de alta complexidade.}
\label{fig:conv2750511b}
\end{figure}

\comoler{este hospital tem perfil de baixo óbito (mortalidade sempre
abaixo de 1\%, típica de maternidade), e por isso sua linha preta corre
bem abaixo das referências. Também aqui há apenas um ano após a
conversão.}

\section{O que anda junto com o quê}
\label{sec:correlacao}

A Figura~\ref{fig:correlacao} resume, num único mapa, o quanto cada par
de indicadores anda junto. A medida usada (correlação de Spearman) vai
de 0, quando dois indicadores não têm relação, a 1, quando crescem
sempre juntos. A mortalidade ajustada foi retirada do mapa por ser
estatisticamente idêntica à mortalidade geral (correlação de 1,000
entre as duas; todas as demais relações coincidem casa a casa).

Os pares mais fortes fazem sentido imediato para a gestão:
\textbf{custo por saída e fração de alta complexidade (0,63)}, custo e
tempo de permanência (0,56), e as duas ocupações entre si (0,48). A
\textbf{complexidade estrutural conversa com tudo}: 0,55 com alta
complexidade, 0,57 com ocupação de internação, 0,49 com custo, 0,47 com
TMP e 0,28 com mortalidade, o retrato numérico do papel de grande
divisor que a Seção~\ref{sec:complexidadeporte} descreveu. Já a
mortalidade se associa mais à ocupação de UTI (0,40) do que à de
internação (0,17), na direção esperada: onde a casuística é mais grave,
usa se mais UTI e morre se mais.

\begin{figure}[H]
\centering
\includegraphics[width=.96\textwidth]{fig_ae_10_correlacao_spearman.png}
\caption{Mapa de correlações entre os indicadores, as ocupações e a
complexidade estrutural. Azul mais escuro indica relação mais forte;
todos os valores estão impressos.}
\label{fig:correlacao}
\end{figure}

\comoler{procure a linha de um indicador e percorra as colunas: cada
célula diz o quanto aquele par anda junto (0 sem relação, 1 relação
perfeita). A diagonal vale sempre 1 porque compara o indicador com ele
mesmo. Repare na coluna da complexidade estrutural, escura quase
inteira: é o confundidor que precisa ser descontado em qualquer
comparação.}

A Figura~\ref{fig:variancia} e a Tabela~\ref{tab:variancia} respondem a
outra pergunta essencial: \textbf{quando um indicador varia, essa
variação vem de comparar hospitais diferentes ou do mesmo hospital
mudando ao longo dos anos?} A resposta é contundente: a variação vem
quase toda das diferenças entre hospitais. No tempo de permanência,
\textbf{apenas 2\% da variação total acontece dentro do mesmo hospital}
ao longo da década; na alta complexidade são 4,4\%, no custo 8,3\%, na
mortalidade 13,9\%, e só a ocupação de UTI passa de um terço (34,7\%).
Em termos de gestão: \textbf{cada hospital é muito parecido consigo
mesmo, ano após ano, e muito diferente dos vizinhos}. Os indicadores
têm forte inércia institucional, e mudanças reais dentro de um hospital
são eventos raros e valiosos de observar, o que dá ainda mais peso aos
cinco casos de conversão da Seção~\ref{sec:conversores}.

\begin{table}[H]
\centering
\caption{De onde vem a variação de cada indicador: parcela que ocorre
dentro do mesmo hospital ao longo dos anos e parcela que vem das
diferenças entre hospitais.}
\label{tab:variancia}
\small
\begin{tabular}{lrr}
\toprule
Indicador & Dentro do hospital (\%) & Entre hospitais (\%) \\
\midrule
TMP                    & 2,0  & 98,0 \\
Fração alta complex.   & 4,4  & 95,6 \\
Custo por saída        & 8,3  & 91,7 \\
Mortalidade geral      & 13,9 & 86,2 \\
Mortalidade ajustada   & 13,9 & 86,2 \\
Ocupação internação    & 17,9 & 82,1 \\
Ocupação UTI           & 34,7 & 65,3 \\
\bottomrule
\end{tabular}
\end{table}

\begin{figure}[H]
\centering
\includegraphics[width=.96\textwidth]{fig_ae_14_variancia_within.png}
\caption{Parcela da variação de cada indicador que acontece dentro do
mesmo hospital ao longo do tempo.}
\label{fig:variancia}
\end{figure}

\comoler{quanto menor a barra, mais estável o indicador dentro de cada
hospital e mais as diferenças observadas são traço permanente da
instituição. O TMP, com 2\%, é quase uma impressão digital do hospital;
a ocupação de UTI, com 35\%, é a que mais responde ao tempo, em boa
parte por causa da pandemia.}

\section{Qualidade dos registros e cuidados adotados}
\label{sec:qualidade}

Nenhuma base administrativa é perfeita, e parte do trabalho
exploratório é localizar exatamente onde estão os problemas para que
eles não contaminem as conclusões. A boa notícia: os problemas deste
painel são poucos, identificados e circunscritos.

\textbf{Ocupação acima de 100\%.} Na internação, 14,7\% das observações
(509) ultrapassam 100\%, e na UTI, 6,1\% (209). Em sua maioria isso é
informação verdadeira de sobrecarga (mais diárias que a capacidade
cadastrada comporta) e não deve ser apagado. O problema real está na
cauda extrema da UTI: \textbf{os 15 maiores valores pertencem todos a
2020 e 2021}, liderados pelo \textbf{Hospital Irmãos Penteado e Santa
Casa de Campinas (CNES 2022648)}, que registra 875\% de ocupação de UTI
em 2021, com 6.389 diárias para apenas 16 leitos cadastrados. A
explicação é operacional: na onda da pandemia os hospitais expandiram
leitos de UTI às pressas e o cadastro não acompanhou, produzindo
divisões por um denominador defasado. O cuidado adotado: analisar a
ocupação como razão positiva, tratar a UTI em duas partes (ter ou não
UTI ativa; e, entre os que têm, a intensidade de uso) e aparar a cauda
extrema (winsorização no percentil 99: 124\% na internação e 164\% na
UTI), com atenção redobrada ao biênio pandêmico.

\textbf{Denominadores frágeis na pandemia.} Ao excluir as internações
COVID, alguns hospitais poderiam ficar com pouquíssimas internações
restantes, tornando suas taxas instáveis. No painel isso ocorre em
\textbf{apenas 2 observações, ambas do Hospital Infante D.\ Henrique
(CNES 2097613, São José do Rio Preto)}: 13 saídas não COVID em 2020 e
somente 3 em 2021, ano em que o TMP calculado degenera para zero (o
único zero de TMP de todo o painel). As duas observações recebem
marcação de denominador frágil e ficam de fora das análises de
resultado.

\textbf{Ruído estatístico dos hospitais pequenos.} A
Figura~\ref{fig:funil} organiza essa discussão: cada ponto é um
hospital num ano, posicionado pelo volume de saídas (eixo horizontal) e
pela mortalidade (eixo vertical); as linhas mostram a faixa em que a
mortalidade deveria variar se as diferenças fossem puro acaso
estatístico, faixa que é larga para hospitais pequenos e estreita para
os grandes. O resultado: \textbf{cerca de 69\% das observações caem
fora até da faixa mais generosa} (1.119 acima e 1.269 abaixo). Isso diz
duas coisas ao mesmo tempo: as diferenças de mortalidade entre
hospitais são reais e estruturais, ligadas ao perfil assistencial, e
não flutuação de sorte; e, abaixo de mil saídas por ano, o acaso passa
a pesar de verdade, o que justifica a cautela extra com hospitais
pequenos.

\textbf{Valores extremos e completude.} Pela regra usual de detecção
(3 intervalos interquartílicos acima do terceiro quartil), há 11
observações extremas de mortalidade (acima de 18,9\%), 82 de custo
(acima de R\$~4.727), 206 de alta complexidade e 223 de TMP, estas
últimas concentradas nos hospitais de longa permanência já descritos,
que merecem tratamento em separado e não corte cego. Quanto à
completude, \textbf{o painel não tem nenhum valor faltante nos sete
indicadores} nem nas variáveis de estratificação; as únicas lacunas
estão em campos cadastrais (nome fantasia e especialização), sem
consequência para as análises.

\begin{figure}[H]
\centering
\includegraphics[width=.96\textwidth]{fig_ae_16_funnel_mortalidade.png}
\caption{Funil da mortalidade: cada ponto é um hospital num ano; a
linha horizontal é a mortalidade média da rede (5,4\%) e as linhas
curvas delimitam a faixa esperada por puro acaso estatístico (95\% e
99,8\%), conforme o volume de saídas.}
\label{fig:funil}
\end{figure}

\comoler{o funil é estreito à direita (hospitais grandes, onde o acaso
pesa pouco) e largo à esquerda (hospitais pequenos, onde uma morte a
mais ou a menos muda a taxa). Pontos fora das linhas curvas indicam
mortalidade que o acaso não explica: a nuvem inteira transbordando das
linhas mostra que as diferenças entre hospitais são estruturais, e a
região à esquerda de mil saídas mostra onde as taxas devem ser lidas
com reserva.}

\section{Como estimamos os efeitos: método e salvaguardas}
\label{sec:estrategia}

Até aqui o documento descreveu os dados. Nesta segunda parte, entramos
na fase de \textbf{estimação}: usamos modelos estatísticos para medir
quanto de cada diferença entre hospitais pode ser atribuído ao modelo
de gestão, depois de descontar complexidade, porte e ano. A ideia
central é simples: comparar hospitais como se estivessem em igualdade
de condições. Um hospital OSS de grande porte e alta complexidade não
pode ser comparado diretamente com uma Santa Casa pequena; o modelo faz
esse desconto de forma sistemática, para todos os hospitais ao mesmo
tempo.

Cada indicador recebeu a família de modelo decidida empiricamente na
primeira parte deste documento: a mortalidade usa uma regressão Beta
com inflação em zero (que acomoda tanto a escala de porcentagens quanto
os raros anos sem óbito); a fração de alta complexidade usa o modelo
ZOIB (que trata separadamente os hospitais sem nenhuma produção de alta
complexidade e os raros casos de produção exclusiva); o custo real por
saída e o tempo de permanência são modelados no logaritmo, comparando
cada hospital com ele mesmo ao longo do tempo (os chamados efeitos
fixos); a produção usa a Binomial Negativa, adequada para contagens com
grande dispersão; a ocupação de internação é modelada no logaritmo sem
truncar em 100\%; e a ocupação de UTI usa um modelo em duas partes,
primeiro se o hospital tem UTI ativa, depois com que intensidade a usa.

Cinco salvaguardas valem para todos os modelos. Primeira: o custo entra
sempre \textbf{em valores reais}, deflacionado pelo IPCA para preços de
2025, porque a primeira parte mostrou que o crescimento nominal do
custo era, na íntegra, inflação. Segunda: as duas observações de
denominador frágil do Hospital Infante D Henrique (CNES 2097613, São
José do Rio Preto, 2020 e 2021, com 13 e 3 saídas fora da COVID) ficam
fora de todas as equações. Terceira: os 18 hospitais de longa
permanência (perfil psiquiátrico e de cuidados prolongados, com
permanência mediana acima de 20 dias) recebem tratamento próprio e
ficam fora das equações de tempo de permanência. Quarta: as ocupações
extremas são limitadas ao percentil 99 (124,2\% na internação e
164,2\% na UTI), preservando a informação de sobrecarga real e
removendo o erro de cadastro documentado do biênio pandêmico. Quinta e
mais importante: \textbf{todas as margens de erro são calculadas
agrupando por hospital}, reconhecendo que 11 anos do mesmo CNES não são
11 informações independentes; sem esse agrupamento, os efeitos
pareceriam muito mais precisos do que de fato são.

Um cuidado adicional merece explicação. Como quase toda a variação dos
indicadores está \emph{entre} hospitais e não \emph{dentro} deles ao
longo do tempo (a Seção~\ref{sec:correlacao} documentou parcelas de
variação interna de 2\% a 35\%), o contraste entre gestões
pode ser lido de duas maneiras: comparando hospitais diferentes
(between) ou acompanhando os cinco hospitais que trocaram de gestão
(within). Na mortalidade usamos a especificação de Mundlak, que separa
as duas leituras dentro do mesmo modelo; no custo e no tempo de
permanência, os efeitos fixos de hospital fazem a leitura within pura.
E a ressalva permanente continua de pé: a categoria Privado tem apenas
3 hospitais e entra nos modelos somente para não contaminar as demais
comparações; nenhum número dela deve ser interpretado.

\section{Resultados por indicador}
\label{sec:resultados}

A Figura~\ref{fig:contrasteoss} resume o contraste central da pesquisa:
quanto cada indicador difere entre a gestão por OSS e a administração
Direta, já descontados complexidade, porte, perfil de longa
permanência e ano, com as margens de erro agrupadas por hospital.

O quadro que emerge é o de um \textbf{regime operacional claramente
mais intenso nas OSS}: produção de saídas \textbf{45,4\% maior},
ocupação de internação \textbf{23,6\% maior} e uso da UTI, entre os
hospitais que a possuem, \textbf{20,4\% maior}, todos com significância
estatística confortável. É uma diferença de regime entre redes com
missões distintas, não um efeito demonstrado da conversão de um
hospital individual (ver as Seções~\ref{sec:efeitogestao}
e~\ref{sec:inferenciafinal}). O \textbf{tempo de permanência é 10,2\%
menor} nas OSS na comparação de cada hospital consigo mesmo (p igual a
0,021, valor que, agrupado por hospital com apenas cinco tratados,
tende a ser otimista; a Seção~\ref{sec:inferenciafinal} o submete a
permutação, que devolve p igual a 0,153), o que fecha a conta do giro:
mais pacientes passam pelos mesmos
leitos porque cada um fica menos tempo internado. O \textbf{custo real
por saída} aparece 6,7\% menor, mas com margem de erro grande demais
para afirmar diferença (o intervalo vai de 22,6\% de queda a 12,4\% de
alta): com apenas cinco hospitais mudando de gestão, os dados não
permitem cravar se as OSS são mais baratas ou mais caras por paciente
atendido. Os efeitos de ano do mesmo modelo confirmam, agora em
ambiente controlado, a queda real de custo da década: 2025 fica 11,6\%
abaixo de 2015 em preços comparáveis.

Na \textbf{mortalidade}, o modelo Beta com inflação em zero e
especificação de Mundlak estima, para a mudança de Direta para OSS
dentro do mesmo hospital, uma redução equivalente a \textbf{1,12 ponto
percentual} na taxa (p igual a 0,085, significância marginal); esse
número não sobrevive intacto ao confronto com o controle sintético da
Seção~\ref{sec:inferenciafinal} e não deve ser citado isoladamente. Duas
leituras complementares importam. Primeiro, o termo de Mundlak (que
captura a diferença estrutural entre hospitais OSS e os demais) é
positivo e não significativo: as OSS não são hospitais de mortalidade
estruturalmente baixa; a queda estimada vem da trajetória dos
conversores. Segundo, o componente de zeros do modelo confirmou um
fato notável da primeira parte: \textbf{nenhum hospital OSS ou
Municipal registra ano sem óbito} no painel, de modo que a chance de
mortalidade zero foi modelada apenas para Filantrópicos e Privados
contra o bloco público. Na \textbf{fração de alta complexidade}, o
modelo ZOIB não encontra diferença entre OSS e Direta em nenhum dos
componentes; quem explica essa fração é a complexidade estrutural (nos
zeros) e o porte (na média), como esperado.

\begin{figure}[H]
\centering
\includegraphics[width=.96\textwidth]{fig_est_03_contraste_oss_direta.png}
\caption{Contraste OSS vs Direta nos modelos principais, em efeito
percentual com intervalo de 95\%; a mortalidade, que tem escala
própria, aparece anotada no título. Custo e TMP: identificação
\emph{within} (efeitos fixos, sustentada pelos cinco conversores);
produção e ocupações: contraste predominantemente \emph{between} entre
as redes OSS e Direta, de natureza associativa.}
\label{fig:contrasteoss}
\end{figure}

\comoler{cada linha é um indicador; o quadrado verde é a melhor
estimativa do efeito de ser OSS em vez de Direta, e a haste horizontal
é a margem de erro. Quando a haste cruza a linha vertical do zero, como
no custo, não se pode afirmar diferença; quando fica inteira de um
lado, como na produção e na ocupação, a diferença é estatisticamente
firme. Valores à direita do zero indicam OSS maior; à esquerda, OSS
menor.}

\section{O efeito da mudança de gestão: estudo de eventos}
\label{sec:efeitogestao}

A forma mais exigente de medir o efeito da gestão OSS é acompanhar
\textbf{os cinco hospitais que de fato mudaram} (Conjunto Hospitalar de
Sorocaba, CNES 2081695; Pérola Byington, CNES 2078287; Hospital Kátia
de Souza Rodrigues de Taipas, CNES 2082225; Hospital Geral de Vila
Penteado, CNES 2091755; e Hospital Estadual de Presidente Prudente,
CNES 2750511) ano a ano em torno da conversão, comparando cada um com
sua própria história e com a evolução simultânea dos demais 309
hospitais. É o chamado \textbf{estudo de eventos}: o ano anterior à
mudança serve de marco zero, e o modelo estima quanto o hospital se
afastou do próprio padrão em cada ano antes e depois da conversão.

As Figuras~\ref{fig:eventomorttmp} e~\ref{fig:eventocustoocup} mostram
os resultados para quatro indicadores. Depois da conversão, a
\textbf{mortalidade cai 1,30 ponto percentual já no ano da mudança},
aprofunda para 2,61 e 2,86 pontos nos dois anos seguintes e estabiliza
em 2,18 pontos na janela de três ou mais anos. O \textbf{tempo de
permanência} cai de forma crescente até 24,6\% na janela longa. O
\textbf{custo real por saída} não reage no ano da mudança, mas acumula
queda de 12,1\% três ou mais anos depois. A \textbf{ocupação de
internação} sobe após a conversão, com pico no primeiro ano.

Antes de comemorar, é preciso olhar o lado esquerdo das figuras, os
anos \emph{anteriores} à conversão. Ali mora o teste de credibilidade:
se o hospital já vinha melhorando antes de mudar de gestão, parte da
melhora posterior não pode ser atribuída à OSS. O teste formal (os
pontos à esquerda da linha tracejada serem todos estatisticamente
nulos) \textbf{passa para o tempo de permanência (p igual a 0,251) e
para a fração de alta complexidade (p igual a 0,344), fica no limite
para o custo (p igual a 0,031) e é rejeitado para a mortalidade e a
ocupação} (p abaixo de 0,001). A rejeição na mortalidade tem nome e
sobrenome, já antecipados na Seção~\ref{sec:conversores}: em Sorocaba
o pico de mortalidade imediatamente anterior à conversão sugere
reversão à média, e no Pérola Byington a melhora começou antes da
mudança formal. Em resumo: \textbf{a evidência mais defensável de
efeito causal da gestão OSS está no tempo de permanência}; para a
mortalidade e a ocupação, os números medem uma associação temporal
real, mas não um efeito causal limpo.

\begin{figure}[H]
\centering
\subfig{fig_est_01_evento_mort_all.png}{Mortalidade geral, em pontos
percentuais}
\subfig{fig_est_01_evento_tmp.png}{Tempo de permanência, em log pontos}
\caption{Estudo de eventos da conversão para OSS: mortalidade geral e
tempo de permanência, amostra completa.}
\label{fig:eventomorttmp}
\end{figure}

\begin{figure}[H]
\centering
\subfig{fig_est_01_evento_custo_real.png}{Custo real por saída, em log
pontos}
\subfig{fig_est_01_evento_ocupacao_internacao_w.png}{Ocupação de
internação, em log pontos}
\caption{Estudo de eventos da conversão para OSS: custo real por saída
e ocupação de internação, amostra completa.}
\label{fig:eventocustoocup}
\end{figure}

\comoler{cada ponto é o desvio do hospital convertido em relação ao seu
próprio padrão, no ano indicado; a haste vertical é a margem de erro e
o losango escuro marca o ano de referência, imediatamente anterior à
mudança. À esquerda da linha tracejada o hospital ainda era Direta; à
direita, já era OSS. Efeito crível pede pontos próximos de zero à
esquerda e afastamento só à direita: é o que se vê no tempo de
permanência, mas não na mortalidade, onde os anos anteriores já
oscilam. Log pontos podem ser lidos, aproximadamente, como variação
percentual dividida por 100.}

\begin{figure}[H]
\centering
\includegraphics[width=.96\textwidth]{fig_est_01_evento_pct_alta_complex.png}
\caption{Estudo de eventos da conversão para OSS: fração de alta
complexidade, em pontos percentuais, amostra completa.}
\label{fig:eventopctalta}
\end{figure}

\comoler{mesma leitura das figuras anteriores. Aqui nenhum ponto
posterior à conversão se firma longe do zero: a mudança de gestão não
alterou o perfil de complexidade da produção desses hospitais.}

Como exercício de robustez, o estudo foi refeito na \textbf{região de
suporte comum} (faixas Barcelona 3 e 4, portes médio e grande), onde
Direta e OSS convivem em número suficiente. Sorocaba (faixa 5) e
Presidente Prudente (faixa 2) saem dessa amostra, que fica com Pérola
Byington, Taipas e Vila Penteado. A Figura~\ref{fig:eventosuporte}
mostra o resultado para mortalidade e tempo de permanência: os efeitos
posteriores ficam ainda maiores (mortalidade 2,06 pontos menor no ano
da conversão, tempo de permanência 18,6\% menor), mas o teste dos anos
anteriores é rejeitado em todos os indicadores e a janela posterior se
apoia em três hospitais, dois deles com um único ano observado após a
mudança. Esses números confirmam a direção dos achados, sem
constituir prova independente.

\begin{figure}[H]
\centering
\subfig{fig_est_02_evento_suporte_mort_all.png}{Mortalidade geral, em
pontos percentuais}
\subfig{fig_est_02_evento_suporte_tmp.png}{Tempo de permanência, em log
pontos}
\caption{Estudo de eventos na região de suporte comum (faixas 3 e 4,
portes médio e grande): mortalidade geral e tempo de permanência.}
\label{fig:eventosuporte}
\end{figure}

\comoler{mesma leitura do estudo de eventos anterior, agora só com
hospitais comparáveis entre si. Note que restam menos anos após a
conversão: dos três conversores presentes, dois mudaram em 2025.}

\section{Robustez e diagnóstico dos modelos}
\label{sec:robustez}

Três baterias de verificação fecham a fase de estimação.

\textbf{Sem o biênio pandêmico.} Todos os modelos foram reestimados
excluindo 2020 e 2021, os anos em que a casuística, os denominadores e
o cadastro de leitos sofreram os choques documentados na primeira
parte. A Figura~\ref{fig:robustezpandemia} compara os contrastes OSS vs
Direta nas duas amostras: as diferenças são mínimas (o tempo de
permanência passa de 10,2\% para 11,3\% de queda, a produção de 45,4\%
para 46,6\% de alta, e assim por diante). \textbf{Nenhuma conclusão
deste documento depende dos anos de pandemia.}

\begin{figure}[H]
\centering
\includegraphics[width=.96\textwidth]{fig_est_05_robustez_sem_pandemia.png}
\caption{Contraste OSS vs Direta estimado no painel completo e na
amostra sem 2020 e 2021, para os modelos em escala percentual.}
\label{fig:robustezpandemia}
\end{figure}

\comoler{cada linha traz duas estimativas do mesmo efeito: o círculo
azul escuro usa todos os anos e o quadrado verde exclui 2020 e 2021.
Quanto mais próximos os dois símbolos, mais robusto o resultado; aqui
eles praticamente se sobrepõem em todas as linhas.}

\textbf{Mortalidade ajustada.} O modelo de mortalidade foi refeito com
a versão que exclui óbitos fetais, maternos e neonatais: o efeito OSS
within vai de 0,230 para 0,233 negativos no logito e o efeito marginal
médio de 1,12 para 1,14 ponto percentual, confirmando que as duas
versões do indicador contam a mesma história, como a correlação
perfeita da primeira parte já anunciava.

\textbf{Resíduos quantílicos.} A primeira parte prometeu que a escolha
das famílias de distribuição seria validada nos modelos ajustados, e é
o que as Figuras~\ref{fig:qqprimeira} e~\ref{fig:qqsegunda} fazem. O
resíduo quantílico é um teste de qualidade: se o modelo descreve bem os
dados, esses resíduos se alinham sobre a diagonal do gráfico. O
alinhamento é muito bom na mortalidade (correlação de 0,984), na
fração de alta complexidade (0,991) e na produção (0,968), e bom no
custo (0,970) e no tempo de permanência (0,952), cujas caudas pesadas
já estavam diagnosticadas e são absorvidas pelas margens de erro
agrupadas. A exceção honesta é a \textbf{intensidade de ocupação da
UTI} (0,814): a cauda de UTIs com uso quase nulo escapa do modelo em
logaritmo, e esse resultado específico deve ser lido com reserva; fica
registrado para a próxima iteração o uso de uma família com massa
perto de zero para esse indicador.

\begin{figure}[H]
\centering
\subfig{fig_est_04_qq_mort_zib.png}{Mortalidade geral, modelo Beta com
inflação em zero}
\subfig{fig_est_04_qq_pct_alta_zoib.png}{Fração de alta complexidade,
modelo ZOIB}
\caption{Resíduos quantílicos dos modelos de taxa.}
\label{fig:qqprimeira}
\end{figure}

\begin{figure}[H]
\centering
\subfig{fig_est_04_qq_producao_bn.png}{Produção, Binomial Negativa}
\subfig{fig_est_04_qq_ocup_uti_pos.png}{Ocupação de UTI entre os
positivos, o único ajuste insatisfatório}
\caption{Resíduos quantílicos da produção e da parte positiva da UTI.}
\label{fig:qqsegunda}
\end{figure}

\comoler{quanto mais os pontos azuis acompanham a linha tracejada,
melhor o modelo descreve os dados; o valor de r no canto é essa
aderência resumida em um número (1 seria perfeito). Compare o
alinhamento quase exato da fração de alta complexidade com o descolamento
da cauda esquerda na UTI, que motivou a ressalva do texto.}

Fecham o quadro as ressalvas permanentes de leitura. A identificação
do efeito de gestão dentro de hospital repousa em \textbf{cinco
conversores}, três deles com um único ano após a mudança; as margens
de erro agrupadas com tão poucas unidades tratadas tendem a ser
otimistas, e os efeitos devem ser lidos como ordens de grandeza. A
categoria Privado (3 CNES) permaneceu como controle de absorção em
todos os modelos, sem interpretação. E os resultados de mortalidade e
ocupação do estudo de eventos carregam a contaminação de tendências
anteriores já descrita, que o tempo de permanência não apresenta.

\section{O teste de estresse final: quanta confiança os efeitos merecem}
\label{sec:inferenciafinal}

A Seção~\ref{sec:robustez} encerrou com um aviso: todo o efeito de
gestão medido dentro de hospital repousa em cinco conversores, e
margens de erro calculadas com tão poucas unidades tratadas tendem a
ser otimistas. Esta seção fecha a pesquisa quantificando exatamente
isso, com três instrumentos usados na literatura para situações de
poucos tratados, e resolve a pendência técnica da UTI.

\textbf{Conversores de mentira.} O teste mais intuitivo é o de
permutação: sorteamos 5 hospitais falsos entre os 309 que nunca
mudaram de gestão, fingimos que eles converteram nos mesmos anos dos
conversores reais e medimos o efeito que o modelo atribuiria a essa
conversão inventada. Repetindo o sorteio 1.999 vezes, obtemos a régua
do acaso. O resultado: a queda de custo real atribuída à OSS é comum
entre os falsos conversores (49\% deles produzem efeito tão grande
quanto o real), enquanto \textbf{a queda de tempo de permanência de
10,2\% é maior que a de 85\% dos placebos} (p igual a 0,153). Um
segundo método, o bootstrap selvagem por cluster, sabidamente
conservador nesta situação, aponta na mesma direção. A tradução: o
efeito no tempo de permanência é consistente em direção e tamanho, mas
\textbf{cinco hospitais não bastam para sustentar a afirmação com o
rigor estatístico convencional}; ele deve ser tratado como evidência
sugestiva forte, não como fato demonstrado.

\textbf{Um estimador moderno para conversões em anos diferentes.} Como
as conversões ocorreram em 2019, 2023 e 2025, a comparação clássica
pode misturar hospitais já convertidos no grupo de controle. O
estimador de Callaway e Sant'Anna evita isso comparando cada coorte de
conversão apenas com os hospitais que nunca mudaram. A
Figura~\ref{fig:cstwfe} mostra que, no ano da conversão, os dois
métodos coincidem em todos os indicadores (na mortalidade, 1,298
contra 1,301 pontos percentuais de queda); nas janelas intermediárias
o método clássico exagera um pouco o efeito do tempo de permanência
(0,25 contra 0,15 em log pontos) e as janelas longas se aproximam, com
intervalos sobrepostos (na mortalidade, 2,18 pontos negativos no
método clássico contra 2,38 no alternativo; no tempo de permanência,
0,28 contra 0,30 em log pontos). Nenhuma conclusão muda; a magnitude
intermediária do tempo
de permanência era o único ponto inflado.

\begin{figure}[H]
\centering
\subfig{fig_est_07_cs_twfe_mort.png}{Mortalidade geral, em pontos
percentuais}
\subfig{fig_est_07_cs_twfe_tmp.png}{Tempo de permanência, em log pontos}
\caption{Estudo de eventos estimado pelo método clássico (TWFE, com
intervalo de confiança) e por Callaway e Sant'Anna, que compara cada
coorte de conversão só com hospitais nunca convertidos.}
\label{fig:cstwfe}
\end{figure}

\comoler{os círculos escuros repetem o estudo de eventos da
Seção~\ref{sec:efeitogestao}; os quadrados verdes são o estimador
alternativo. Quando os dois símbolos quase se tocam, como no ano zero
e na janela longa, o resultado não depende do método; a distância nos
anos intermediários do tempo de permanência mostra o único ponto em
que o método clássico exagerava.}

\textbf{O gêmeo sintético.} Para Sorocaba e Pérola Byington, os dois
conversores com janela posterior utilizável, construímos um
\emph{controle sintético}: uma combinação ponderada de hospitais nunca
convertidos que reproduz a trajetória do hospital antes da mudança de
gestão. Se a conversão teve efeito, o hospital real deve descolar do
seu gêmeo depois dela. A Figura~\ref{fig:sintetico} traz o resultado
para a mortalidade, e ele é a ducha de água fria mais importante do
documento: \textbf{em Sorocaba o hospital real fica 1,09 ponto
percentual acima do gêmeo} após a conversão (o gêmeo, que acompanhava
a subida até 2018, também cai depois, confirmando que parte da queda
era retorno natural ao patamar); \textbf{no Pérola Byington o hospital
fica 0,95 ponto abaixo do gêmeo}, na direção esperada, mas dentro do
que sorteios de placebo produzem (p igual a 0,33). No custo real, a
lacuna do Pérola é de R\$ 354 a menos por saída (p igual a 0,28). Em
nenhum caso o descolamento é raro o bastante para valer como prova.

\begin{figure}[H]
\centering
\subfig{fig_est_06_sintetico_2081695_mort_pp.png}{Conjunto Hospitalar
de Sorocaba (CNES 2081695), conversão em 2019}
\subfig{fig_est_06_sintetico_2078287_mort_pp.png}{Pérola Byington
(CNES 2078287), conversão em 2023}
\caption{Mortalidade geral: hospital real contra o seu controle
sintético, construído só com hospitais nunca convertidos.}
\label{fig:sintetico}
\end{figure}

\comoler{a linha preta é o hospital; a azul tracejada é o gêmeo
sintético, montado para imitar o hospital antes da linha vertical. O
efeito da conversão é a distância entre as duas linhas depois da linha
vertical. Em Sorocaba a linha preta termina acima da azul; no Pérola,
abaixo, mas com o gêmeo também caindo.}

\textbf{A pendência da UTI, resolvida.} A intensidade de uso de UTI
foi reestimada com a família Gama, adequada à cauda de UTIs de uso
quase nulo que o modelo em logaritmo não acomodava. O diagnóstico de
resíduos salta de 0,814 para \textbf{0,944}
(Figura~\ref{fig:qqgama}, no Apêndice~\ref{ap:diagnosticos}),
e o resultado substantivo permanece: uso de UTI \textbf{18,2\% maior
nas OSS} (p igual a 0,005), com a variante de corte mínimo de
atividade confirmando (22,3\%). A conclusão da
Seção~\ref{sec:resultados} fica, portanto, validada também neste
indicador.

\textbf{Balanço final da fase de estimação.} O quadro que resiste a
toda a bateria de testes é: as OSS operam um regime assistencial
claramente mais intenso (mais produção, mais ocupação, mais uso de
UTI), com diferenças firmes \textbf{enquanto retrato transversal das
redes}; a única evidência com pretensão causal que resiste é a do
tempo de permanência, presente em todos os métodos, e mesmo ela
permanece sugestiva com cinco conversores, não definitiva; o custo
real por saída não mostra diferença demonstrável em nenhum método; e
a queda de mortalidade após as conversões, embora visível, não
sobrevive ao confronto com o gêmeo sintético em Sorocaba e fica sem
significância no Pérola Byington. A segunda fase da estimação
(Seção~\ref{sec:fase2}) quantificou esse último ponto da forma mais
direta possível: retirando os cinco conversores da amostra, a
diferença de mortalidade entre as redes OSS e Direta simplesmente
desaparece (o coeficiente vira 0,09 positivo, estatisticamente nulo),
ou seja, \textbf{a mortalidade menor é um fenômeno da trajetória dos
hospitais que mudaram de gestão, sem contraparte na comparação entre
as redes}. Novos anos de dados, especialmente
para os três hospitais convertidos em 2025, são o caminho natural para
transformar as sugestões em conclusões.

\section{A segunda camada de evidência: hierarquia e eficiência}
\label{sec:fase2}

A segunda fase da estimação atacou a mesma pergunta com três famílias
de métodos ainda não usadas no documento: modelos
\textbf{hierárquicos} (que tratam os 11 anos de cada hospital como
parentes, não como observações soltas), a análise envoltória de dados
(\textbf{DEA}, que constrói uma fronteira de melhores práticas e mede
a distância de cada hospital até ela) e a fronteira estocástica
(\textbf{SFA}, a contraparte paramétrica do DEA). O pacote completo de
resultados, com scripts, diagnósticos e um relatório de execução
auditável, está na pasta resultados\_fase2 do repositório; aqui ficam
os achados que mudam ou reforçam a leitura do documento.

\textbf{O que a camada hierárquica confirma.} Reestimados com um
efeito próprio para cada hospital, o tempo de permanência e o custo
real reproduzem quase exatamente os resultados da
Seção~\ref{sec:resultados} (queda de 10,2\% no TMP; custo sem
diferença), a produção e a ocupação continuam maiores nas OSS com
magnitudes menores (33\% e 13\%, porque parte do contraste entre as
redes é absorvida pelo efeito de hospital), e a mortalidade aparece
0,255 menor no logito, com precisão maior que a da
Seção~\ref{sec:resultados}. O teste decisivo, porém, é outro:
\textbf{retirando os cinco conversores, a diferença de mortalidade
entre as redes desaparece} (0,09 positivo, nulo), enquanto TMP,
produção, ocupação e eficiência sobrevivem intactos. É a quantificação
do aviso da Seção~\ref{sec:inferenciafinal}: a mortalidade menor é uma
história dos hospitais que mudaram de gestão, não uma diferença entre
as redes OSS e Direta.

\begin{figure}[H]
\centering
\includegraphics[width=.9\textwidth]{figB02_encolhimento_categorias.png}
\caption{Efeitos de categoria na mortalidade com e sem partial
pooling: com apenas 3 hospitais, o efeito aparente do Privado encolhe
em quase 90\% quando o modelo pondera a evidência pelo tamanho do
grupo.}
\label{fig:encolhimento}
\end{figure}

\comoler{cada linha compara duas estimativas do efeito da mesma
categoria: o círculo escuro trata cada categoria isoladamente; o
quadrado verde deixa as categorias emprestarem informação umas às
outras, na proporção da quantidade de dados de cada uma. Quanto menor
o grupo, maior o recuo do quadrado em direção ao zero: o Privado, com
3 hospitais, quase não sustenta efeito próprio. É a ilustração formal
da regra deste documento de nunca interpretar essa categoria.}

\textbf{Eficiência: a ressalva vem antes do resultado.} O DEA e a SFA
medem eficiência produtiva: quanto produto (saídas, incluída a alta
complexidade) cada hospital extrai dos seus insumos. Um dos insumos é
o \textbf{valor total real} pago pelo SUS, e ele exige a ressalva de
primeira ordem desta seção: o valor do SIH é reembolso por produção
faturada, não medida direta de recursos consumidos, e os regimes de
pagamento diferem por categoria (as OSS recebem pelo contrato de
gestão, de modo que o faturamento por AIH pode divergir do recurso
efetivo; os filantrópicos convivem com subfinanciamento e subsídio
cruzado). Parte da distância entre categorias pode refletir regime de
pagamento, não gestão. Duas verificações sustentam que a ordenação
não é artefato do insumo financeiro: a \textbf{produção física por
leito}, que não usa dinheiro nenhum, já ordena as redes do mesmo jeito
(54 saídas por leito nas OSS e 51 nos Municipais, contra 35 na Direta
e 34 nos Filantrópicos); e o DEA reestimado \textbf{só com leitos como
insumo} preserva e até amplia a vantagem (medianas de 0,68 nas OSS
contra 0,40 na Direta).

Com a ressalva posta, o resultado: na eficiência corrigida de viés
(bootstrap de Simar e Wilson, 2.000 réplicas por ano), as medianas do
período são \textbf{OSS 0,72, Público Municipal 0,64, Direta 0,59 e
Filantrópico 0,58}. Na regressão da distância à fronteira sobre
categoria, complexidade, porte e ano, a vantagem OSS equivale a cerca
de \textbf{40\% menos distância à fronteira que a Direta}, estável
retirando o biênio pandêmico, os cinco conversores ou restringindo ao
suporte comum. Ao contrário da mortalidade, \textbf{a eficiência
maior das OSS é um traço da rede}, presente na comparação pura entre
hospitais distintos. E o receio de os 3 hospitais Privados
distorcerem a fronteira não se confirma: retirados, os escores dos
demais praticamente não se movem (correlação de ordem 1,000), porque
o perfil deles opera longe da fronteira.

\begin{figure}[H]
\centering
\subfig{figE01_violino_te_categoria.png}{Eficiência DEA corrigida de
viés}
\subfig{figS01_te_categoria.png}{Eficiência técnica da fronteira
estocástica}
\caption{Eficiência por categoria nos dois métodos. No painel (b), a
concentração das OSS junto ao valor 1 é uma degeneração numérica do
modelo (o texto explica) e o valor delas não deve ser citado; a ordem
das demais categorias é a informação válida.}
\label{fig:eficiencia}
\end{figure}

\comoler{como nos violinos das seções anteriores: a silhueta engorda
onde há mais hospitais e o número é a mediana do grupo. Compare a
posição dos violinos de OSS e Municipal com os de Direta e
Filantrópico nos dois painéis: métodos de naturezas opostas ordenam as
redes do mesmo jeito. O Privado aparece só como registro.}

\textbf{A contraparte paramétrica e a convergência entre métodos.} A
fronteira estocástica (Translog, preferida ao Cobb Douglas por larga
margem; 94\% da variância composta atribuída à ineficiência) confirma
a ordem das redes, com uma honestidade a registrar: as OSS ficam tão
coladas à fronteira paramétrica que o coeficiente delas na equação de
ineficiência perde identificação (o análogo da separação perfeita da
Seção~\ref{sec:estrategia}), e por isso o NÍVEL da eficiência das OSS
na SFA não é citável; vale a ordem, não o número. O resultado síntese
da fase é a concordância entre os dois métodos: correlação de ordem
de \textbf{0,63 entre as eficiências do DEA e da SFA} (0,65 excluindo
as OSS, mostrando que a concordância não depende do grupo problemático),
crescendo ao longo do período.

\begin{figure}[H]
\centering
\includegraphics[width=.9\textwidth]{figS02_dispersao_sfa_dea.png}
\caption{Convergência entre os dois métodos de eficiência: cada ponto
é um hospital num ano, posicionado pelo escore DEA (horizontal) e pela
eficiência SFA (vertical).}
\label{fig:convergencia}
\end{figure}

\comoler{quanto mais a nuvem sobe da esquerda para a direita, mais os
dois métodos concordam sobre quem é eficiente. A faixa verde colada no
topo é a degeneração das OSS descrita no texto (o modelo paramétrico as
coloca todas na fronteira); mesmo descontando essa faixa, a
concordância entre os métodos permanece.}

\textbf{O que fica para a leitura do documento.} A segunda fase não
altera nenhuma conclusão; ela redistribui os pesos. O \textbf{regime
operacional mais intenso e a eficiência produtiva maior das OSS são
traços robustos da rede}, visíveis em todos os métodos e em todas as
amostras, inclusive sem os conversores; o \textbf{TMP menor} replica
com exatidão em especificação hierárquica; o \textbf{custo} segue sem
diferença em qualquer método; e a \textbf{mortalidade menor}
concentra se inteiramente na trajetória dos cinco conversores, o
que a aproxima das ressalvas causais das
Seções~\ref{sec:efeitogestao} e~\ref{sec:inferenciafinal} e a afasta
de uma leitura de superioridade da rede OSS nesse desfecho.


\appendix
\section{Diagnósticos de ajuste: gráficos QQ e resíduos quantílicos}
\label{ap:diagnosticos}

\setcounter{figure}{0}
\renewcommand{\thefigure}{A.\arabic{figure}}

Este apêndice reúne os gráficos de diagnóstico técnico do documento,
para que o corpo do texto flua sem interrupções; cada figura mantém o
seu bloco de leitura. Duas noções bastam para acompanhar tudo.
\textbf{Gráfico QQ}: compara os dados reais com um molde matemático de
distribuição; se os pontos seguem a reta, o molde descreve bem os
dados, e desvios sistemáticos da reta indicam onde o molde falha. O
número r que acompanha cada QQ resume essa aderência: quanto mais
próximo de 1, melhor. \textbf{Resíduo quantílico}: é a versão do QQ
para um modelo já ajustado com covariáveis; se o modelo descreve bem
os dados, esses resíduos se comportam como uma Normal padrão e os
pontos deitam na diagonal.

\subsection*{Diagnósticos da fase exploratória}

As Figuras~\ref{fig:qqcusto} e~\ref{fig:qqtmp} avaliam os moldes
marginais de faturamento e TMP da Seção~\ref{sec:distribuicoes}.

\begin{figure}[H]
\centering
\subfig{fig_ae_03_qq_lognormal_custo_saida.png}{Molde LogNormal}
\subfig{fig_ae_03_qq_fisk_custo_saida.png}{Molde log logístico (Fisk)}
\caption{Gráficos QQ do faturamento por saída: os pontos são os dados e a
reta é o comportamento esperado pelo molde.}
\label{fig:qqcusto}
\end{figure}

\comoler{quanto mais os pontos abraçam a reta, melhor o molde. No
painel (b), o molde log logístico segue os dados por quase toda a
extensão, descolando apenas nos pouquíssimos hospitais mais caros; por
isso ele recebe a melhor nota da Tabela~\ref{tab:familias}.}

\begin{figure}[H]
\centering
\subfig{fig_ae_03_qq_lognormal_tmp.png}{Molde LogNormal}
\subfig{fig_ae_03_qq_fisk_tmp.png}{Molde log logístico (Fisk)}
\caption{Gráficos QQ do TMP no painel completo.}
\label{fig:qqtmp}
\end{figure}

\comoler{nos dois painéis os pontos acompanham a reta em quase toda a
extensão e descolam apenas nos pouquíssimos hospitais de permanência
mais longa. O painel corrente já não contém a subpopulação de longa
permanência que, na versão anterior desta análise, criava um teto de
30,5 dias e o característico achatamento no alto: a ETAPA F a removeu, e
o ajuste do TMP passou a ser quase perfeito.}

As Figuras~\ref{fig:qqcondA} a~\ref{fig:qqcondC} sustentam a conclusão
da Seção~\ref{sec:distribuicoes} de que o desajuste dos QQ marginais
vem de misturar 289 hospitais heterogêneos, não do molde.

\begin{figure}[H]
\centering
\subfig{fig_ae_11_qq_custo_marginal.png}{Faturamento: todos os hospitais
juntos}
\subfig{fig_ae_11_qq_custo_residuo_ef.png}{Faturamento: cada hospital
comparado consigo mesmo}
\caption{Faturamento por saída: gráfico QQ antes e depois de controlar as
diferenças fixas entre hospitais e entre anos.}
\label{fig:qqcondA}
\end{figure}

\comoler{compare o gráfico de cima com o de baixo: em cima, todos os
hospitais juntos, os pontos serpenteiam em volta da reta; embaixo,
comparando cada hospital consigo mesmo, os pontos praticamente deitam
na reta.}

\begin{figure}[H]
\centering
\subfig{fig_ae_11_qq_tmp_marginal.png}{TMP: todos os hospitais juntos}
\subfig{fig_ae_11_qq_tmp_marginal_sem_longa_perm.png}{TMP: sem a antiga
subpopulação de longa permanência}
\caption{TMP: gráfico QQ do painel completo e da versão sem a antiga
subpopulação de longa permanência --- no painel corrente as duas
coincidem, porque a ETAPA F já removeu essa subpopulação.}
\label{fig:qqcondB}
\end{figure}

\comoler{no painel anterior este par contrastava o desenho em degrau
(com os hospitais de longa permanência) com a reta quase perfeita (sem
eles). No painel corrente a ETAPA F já removeu essa subpopulação, então
as duas versões coincidem e ambas mostram a reta --- a confirmação de
que o desvio era a subpopulação, não o molde.}

\begin{figure}[H]
\centering
\subfig{fig_ae_11_qq_mort_logito_marginal.png}{Mortalidade: todos os
hospitais juntos}
\subfig{fig_ae_11_qq_mort_logito_residuo_ef.png}{Mortalidade: cada
hospital comparado consigo mesmo}
\caption{Mortalidade: gráfico QQ antes e depois de controlar as
diferenças fixas entre hospitais e entre anos.}
\label{fig:qqcondC}
\end{figure}

\comoler{mesmo padrão do faturamento: a curva em S do gráfico de cima, com
todos os hospitais misturados, vira uma reta bem comportada quando cada
hospital é comparado consigo mesmo. A heterogeneidade entre hospitais,
e não o molde, era a fonte do desajuste.}

\subsection*{Resíduos quantílicos dos modelos estimados}

As Figuras~\ref{fig:qqprimeira} e~\ref{fig:qqsegunda} validam as
famílias dos modelos da Seção~\ref{sec:resultados}, e a
Figura~\ref{fig:qqgama} documenta a reestimação da UTI comentada na
Seção~\ref{sec:inferenciafinal}.

\begin{figure}[H]
\centering
\subfig{fig_est_04_qq_mort_zib.png}{Mortalidade geral, modelo Beta com
inflação em zero}
\subfig{fig_est_04_qq_pct_alta_zoib.png}{Fração de alta complexidade,
modelo ZOIB}
\caption{Resíduos quantílicos dos modelos de taxa.}
\label{fig:qqprimeira}
\end{figure}

\begin{figure}[H]
\centering
\subfig{fig_est_04_qq_producao_bn.png}{Produção, Binomial Negativa}
\subfig{fig_est_04_qq_ocup_uti_pos.png}{Ocupação de UTI entre os
positivos, o único ajuste insatisfatório}
\caption{Resíduos quantílicos da produção e da parte positiva da UTI.}
\label{fig:qqsegunda}
\end{figure}

\comoler{quanto mais os pontos azuis acompanham a linha tracejada,
melhor o modelo descreve os dados; o valor de r no canto é essa
aderência resumida em um número (1 seria perfeito). Compare o
alinhamento quase exato da fração de alta complexidade com o
descolamento da cauda esquerda na UTI, que motivou a reestimação da
figura seguinte.}

\begin{figure}[H]
\centering
\includegraphics[width=.72\textwidth]{fig_est_08_qq_uti_gama.png}
\caption{Resíduos quantílicos da intensidade de UTI reestimada com a
família Gama: a aderência sobe de 0,814 para 0,944.}
\label{fig:qqgama}
\end{figure}

\comoler{compare com o painel (b) da figura anterior: os pontos que
antes despencavam da linha tracejada na ponta esquerda agora a
acompanham em quase todo o trajeto.}


\end{document}
